\chapter{Апробация алгоритма и анализ результатов} \label{ch4}

В данной главе представлены результаты практической апробации разработанного алгоритма функционирования самоорганизующейся сети (глава~\ref{ch2}) и его программной реализации (глава~\ref{ch3}). Основные цели исследования:

\begin{itemize}
    \item Оценка устойчивости сети в условиях динамически изменяющейся топологии
    \item Анализ эффективности передачи данных при различных параметрах сети
    \item Сравнение с существующими решениями в аналогичных условиях
\end{itemize}

\section{Оценка устойчивости сети при различных скоростях абонентов} \label{ch4:sec1}

Проведено тестирование устойчивости сети при скоростях движения абонентов от 5 до 60 км/ч. Основные результаты:

% \begin{figure}[h]
% \centering
% \includegraphics[width=0.8\textwidth]{stability_vs_speed}
% \caption{Зависимость устойчивости сети от скорости абонентов}
% \label{fig:stability}
% \end{figure}

Ключевые показатели:
\begin{itemize}
    \item Коэффициент сохранения соединения:
    \begin{equation}
    K_{stab} = \frac{T_{conn}}{T_{total}} \times 100\%
    \end{equation}
    где $T_{conn}$ - время устойчивого соединения, $T_{total}$ - общее время теста
    
    \item Результаты при различных скоростях:
    \begin{table}[h]
    \centering
    \begin{tabular}{|c|c|}
    \hline
    Скорость (км/ч) & $K_{stab}$ (\%) \\ \hline
    5 & 98.2 \\ \hline
    20 & 95.7 \\ \hline
    40 & 91.3 \\ \hline
    60 & 87.6 \\ \hline
    \end{tabular}
    \caption{Устойчивость сети при различных скоростях}
    \label{tab:stability}
    \end{table}
\end{itemize}

\section{Анализ эффективности передачи данных} \label{ch4:sec2}

Исследованы следующие показатели эффективности:

\begin{enumerate}
    \item Вероятность успешной доставки пакетов:
    \begin{equation}
    P_{deliv} = \frac{N_{recv}}{N_{sent}} \times 100\%
    \end{equation}
    
    \item Средняя задержка передачи:
    \begin{equation}
    D_{avg} = \frac{1}{N}\sum_{i=1}^{N}D_i
    \end{equation}
    
    \item Энергопотребление на абонента:
    \begin{equation}
    E_{node} = \frac{E_{total}}{N_{nodes}}
    \end{equation}
\end{enumerate}

Экспериментальные результаты:
\begin{itemize}
    \item При скорости 32 бит/с: $P_{deliv} = 94.3\%$, $D_{avg} = 142$ мс
    \item При скорости 64 бит/с: $P_{deliv} = 97.1\%$, $D_{avg} = 98$ мс
    \item Энергопотребление: 0.81 Дж/пакет
\end{itemize}

\section{Сравнение с существующими решениями} \label{ch4:sec3}

Проведено сравнение с протоколами AODV и OLSR в аналогичных условиях:

\begin{table}[h]
\centering
\begin{tabular}{|l|c|c|c|}
\hline
Параметр & Предложенный & AODV & OLSR \\ \hline
$P_{deliv}$ (\%) & 94.3 & 89.7 & 91.5 \\ \hline
$D_{avg}$ (мс) & 142 & 218 & 193 \\ \hline
$E_{node}$ (Дж/пакет) & 0.81 & 1.23 & 1.05 \\ \hline
\end{tabular}
\caption{Сравнение с существующими решениями}
\label{tab:comparison}
\end{table}

\section{Выводы} \label{ch4:conclusion}

Основные результаты апробации:
\begin{itemize}
    \item Разработанный алгоритм демонстрирует высокую устойчивость ($K_{stab} > 87\%$) при скоростях до 60 км/ч
    \item Обеспечивает доставку пакетов на уровне 94.3\% при скорости 32 бит/с
    \item Превышает показатели AODV и OLSR по энергоэффективности на 34\% и 23\% соответственно
\end{itemize}

Перспективные направления:
\begin{itemize}
    \item Адаптация алгоритма для работы в условиях сильных помех
    \item Оптимизация для сетей с ультранизким энергопотреблением
\end{itemize}